%======================================================================
%  Supplementary Methods -- model interpretation (paper section 5)
%
%  Source artifact: artifact_interpretation/page.html, section#methods,
%  snapshot 2026-08-28. Seven subsections, transcribed verbatim.
%  The five ".eq" divs are rendered here as real display math; nothing else
%  in the wording was changed. The method-writeups/*.md files in that tree
%  are stale and contradict this text -- they were not used.
%======================================================================

\subsection*{Masked-logit inference}

For a given plasmid, per-position predictive distributions were obtained by
masking each position in turn and recording the model's output logits over
the nucleotide vocabulary, using a 2,048 bp context window centred on the
masked position. Windows are extracted circularly for replicons annotated as
circular; none of the windows displayed here in fact crosses the origin of
its deposited sequence. Because the encoder is not explicitly strand-symmetric,
logits were computed independently for the forward sequence and for its
reverse complement, with reverse-complement predictions re-mapped to
forward-strand coordinates, and the two were averaged in logit space before
the softmax.

\subsection*{Positional information content}

To summarize how sharply the model's predictive distribution is concentrated
at each position, logits were converted to probabilities by softmax after
forward/reverse-complement averaging, and information content was computed as
the Kullback--Leibler divergence of the predicted distribution from a uniform
background over the four bases,

\[
\mathrm{IC}(i) = \sum_b p_{i,b}\,\log_2\!\left(\frac{p_{i,b}}{0.25}\right)
\;\in\; [0,2]\ \text{bits}.
\]

High-IC positions are those at which the model's prediction is sharply
concentrated on one base given its context; low-IC positions are
near-uniform. Information content is a property of the model's output
relative to the stated background and is not a measure of evolutionary
conservation or of biological importance. Per-base contributions
$p_{i,b}\cdot\mathrm{IC}(i)$ set letter heights in the logo tracks.

\subsection*{Functional annotation of plasmids}

Structural and functional annotations for each plasmid were generated with
PlasAnn,
% TODO cite: PlasAnn -- the entry islam2026plasann has been added to
% songlab.bib from the artifact's ref.bib, but the artifact carries no
% citation marker here. Insert \cite{islam2026plasann} if wanted.
run on the assembled nucleotide sequence for each accession and read from
cached per-accession outputs. For the projections shown here the annotation
track was built as follows. PlasAnn's own origin-of-transfer and
origin-of-replication calls were removed and replaced outright by curated
BLAST-derived \emph{oriT} and \emph{oriV} hits for that accession (45 and 85
hits respectively across the cohort), with no further quality filtering at
this step. Mobile-element and replicon-associated calls were discarded rather
than demoted, so their footprint becomes intergenic only where no surviving
feature covers it. Remaining features were mapped onto consecutive 10 bp
tiles, a feature claiming every tile its span touches; a tile touched by two
or more distinct feature types was labelled ambiguous and removed before the
projection was fitted, and tiles touched by nothing were labelled intergenic.
Ambiguity is therefore resolved by exclusion, not by a priority ranking: a
priority hierarchy exists in the codebase but is bypassed on this path.
Interval operations use bioframe.
% TODO cite: bioframe.
Annotations were used after embedding generation, for tile filtering,
colouring and interpretation, and at no point as supervision.

\subsection*{Locus curation and coordinate verification}

All intervals in this section are 0-based half-open; figure axes are labelled
in 1-based genomic coordinates, so an axis reads one greater than the
interval start. The two main-figure loci were taken from their primary
literature --- the pSC101 minimal origin from BAL 31 deletion analysis, and
the pIGRK \emph{oriT1} region from deletion and conjugation assays ---
% TODO cite: the pSC101 and pIGRK primary sources.
and every feature was then located by exact literal substring search against
the downloaded FASTA sequence rather than transcribed from a published
figure; each motif matched exactly once.

Supplementary loci come from mixed sources, and the class of evidence differs
between them. Features at p838B-R, pGNB2 and the pKLM array are taken from
the deposited GenBank curation; features at pHTbeta, NAH7 and R388 are from
primary literature, with coordinates relocated against the deposited
sequence.
% TODO cite: the pHTbeta, NAH7, R388, pTA1040 and PA83 locus sources.
Four departures are recorded explicitly. The pECTm80 direct repeats and
AT-rich interval are heuristic computational annotations from a consensus
scan and a local base-composition measurement, not experimentally mapped
features. The pGNB2 fourth iteron is drawn at 22 bp, an analytical extension
of a 6 bp deposited call to the tandem period of the other three, which it
matches at 77\% identity. The NAH7 reading frame was selected to maximize arm
self-complementarity before plotting, so observations at that locus are
conditional on that choice. The pTA1040 stem-loops are literature dyad
symmetries relocated by minimum-edit-distance search, with edit distances of
1 to 5 over 18 bp; four of the source's five loops are shown, the omitted one
being that the source describes as the most constrained.

The $\sigma^{70}$ promoter panels are selected examples, not a benchmark.
Three (\emph{cop}, \emph{ycbr}, \emph{sacaa}) are hand-picked from a
genome-wide motif scan; their drawn $-$35, $-$10 and Shine--Dalgarno boxes
are scan hits, while the box-pair coupling boxes and the coding-sequence box
are annotation overlays. The pHTbeta ORF32 panel is not from that scan: its
boxes were placed by exact substring search for the canonical consensus
upstream of the start codon, and are therefore expected-position overlays.
CRISPR arrays were called with CRISPRCasFinder
% TODO cite: CRISPRCasFinder.
and only its direct-repeat calls are drawn --- no spacer glyphs and no
model-detected array extension. The caller's 1-based-inclusive Start/End were
converted as (Start$-$1, End) and the conversion verified against
JX524189.1's own GenBank repeat features, which match on all sixteen repeats.
The four displayed arrays are selected from six that were rendered; the
selection was made partly on the basis of the model output, and the omitted
cases are identified in the Supplementary legend.

\subsection*{Categorical Jacobian}

To measure the pairwise sequence dependencies the frozen model carries within
a window, we computed a categorical Jacobian and reduced it to a dependency
map, following the construction used to extract co-evolutionary signal from
protein and genomic language models.
% TODO cite: the categorical-Jacobian construction.
The quantity is a finite difference over discrete substitutions, not a
continuous derivative, and not a measurement of physical contact or
evolutionary covariation. Let $x \in \{\mathrm{A,C,G,T}\}^{L}$ be the
reference sequence of a fixed $L = 2{,}048$ bp context, let $x^{i \to a}$
denote $x$ with position $i$ replaced by allele $a$, and let
$f_{j,b}(\cdot)$ be the model's output logit for allele $b$ at position $j$.
Substituting every position with every allele gives the four-dimensional
tensor

\[
J_{i,a,j,b} = f_{j,b}\!\left(x^{i \to a}\right) - f_{j,b}(x),
\qquad i,j \in [L];\;\; a,b \in \{\mathrm{A,C,G,T}\}.
\]

The tensor is then centred along each of its four axes in turn,

\[
\tilde{J} = (1 - M_i)(1 - M_a)(1 - M_j)(1 - M_b)\,J,
\]

where $M_d$ replaces the tensor by its mean along axis $d$; the four
operators commute, so the order is immaterial. Centring along $b$ subtracts
the position-$j$ log-normalizer, so each entry of $\tilde{J}$ is a log
probability ratio of the form
$\log\!\left[P(n_j = b \mid n_i = a) / P(n_j = b \mid n_i = x_i)\right]$ up to
a constant in $b$, and centring along $a$ removes the corresponding reference
offset. The $4 \times 4$ allele block at each pair of positions is reduced to
a scalar dependency score by its Frobenius norm, which weights every allele
pair rather than any single one:

\[
C_{ij} = \bigl\lVert \tilde{J}_{i,\cdot,j,\cdot} \bigr\rVert_F
= \Bigl(\textstyle\sum_{a,b} \tilde{J}^{\,2}_{i,a,j,b}\Bigr)^{1/2}.
\]

The matrix is symmetrized and its diagonal zeroed, and an average-product
correction is then applied, which removes the background contributed by
generically high-signal positions:

\[
S = \tfrac{1}{2}\!\left(C + C^{\mathsf{T}}\right), \quad S_{ii} = 0;
\qquad
S^{\mathrm{APC}}_{ij} = S_{ij} - \frac{S_{i\cdot}\,S_{\cdot j}}{S_{\cdot\cdot}}.
\]

where $S_{i\cdot}$ and $S_{\cdot j}$ are the row and column sums of $S$ and
$S_{\cdot\cdot}$ its grand sum. The diagonal is zeroed before this correction
and not restored afterwards, so displayed diagonal entries carry the
correction term alone. No positions are masked in this computation: each
perturbed input carries a genuine alternative nucleotide at position $i$ and
the full observed sequence elsewhere. Unlike the per-position logit tracks
described above, the Jacobian is computed from a single forward-orientation
pass and is not averaged with the reverse complement, and no
reverse-complement concordance control is shown. An alternative reduction,
taking the largest signed entry of each allele block, exists in the same code
path but was not used for any panel shown here.

Maps were aligned to genomic coordinates and cropped to the locus for
display, with the annotation and logo tracks drawn in the same coordinate
frame. Each panel is scaled to its own crop, with the lower limit at the crop
minimum and the upper limit at its 99th percentile rather than its maximum.
Two consequences follow and are relied on throughout: absolute intensities
are not comparable between panels, and the diverging colour map is not
centred on zero, so a mid-scale colour does not indicate the absence of
dependency.

\subsection*{Contextual sequence embeddings and dimensionality reduction}

Each plasmid was divided into consecutive, non-overlapping 10 bp tiles. For
each tile, the unmasked 2,048 bp context centred on it was passed through the
frozen model, the final encoder layer's hidden states at the tile's own ten
centre positions were averaged, and the forward-strand and reverse-complement
results were averaged in turn, giving one embedding per tile. The cohort is
40 Enterobacterales plasmids selected for curated \emph{oriT} and \emph{oriV}
BLAST hits and genus-weighted for diversity; it is origin-enriched by
construction and is not a representative sample of plasmids. This yielded
310,013 tiles, of which 1,535 with ambiguous annotation were removed, leaving
308,478 displayed.

Embeddings were standardized across all tiles, then L2-normalized so that a
cosine metric could be used, and projected to two dimensions with UMAP
% TODO cite: UMAP (mcinnes2018umap exists in songlab.bib) and torchdr.
(\texttt{torchdr}; 1,300 neighbours, minimum distance 0.5, 2,000 iterations,
PCA initialization, seed 42). The projection was fitted on the 279,480 unique
embedding rows, collapsing 28,998 exact duplicates, and coordinates were then
broadcast back to every tile. Points were coloured by the annotation track
described above, by a finer annotation vocabulary, or by GC content measured
directly on the same 10 bp tile; with a ten-base denominator, GC takes only
eleven distinct values, and no binning step is applied. Base composition was
necessarily present in the model's input, so GC is not an external property
withheld from the model. These projections use the final encoder layer,
whereas the origin classifier described in \secref{sec:origin} is fitted on
% NOTE (2026-09-01): the artifact read "whereas the §3 origin classifier", and
% this line carried a TODO because \secref renders the heading TEXT. Recast on
% the pattern settled in the §2 split, so the heading is the object of "in".
% Matches the identical recast in 05_interpretation.tex.
\texttt{blocks.60}, twelve residual blocks earlier; the two are different
representations of the same frozen model. No separation metric or classifier
was computed on the projection, so we report association rather than
separation.

\subsection*{Regulatory motif scanning}

Putative $\sigma$-factor promoter elements were identified by scanning plasmid
sequences against position-weight matrices for $\sigma^{70}$ $-$35 and $-$10
elements and for Shine--Dalgarno sequences, derived from RegulonDB, using
FIMO as implemented in pymemesuite.
% TODO cite: RegulonDB, FIMO/MEME Suite, pymemesuite.
The matrix files carry no background frequencies, so scoring fell back to a
uniform nucleotide background; no per-sequence background was estimated and
no multiple-testing correction was applied, so reported p-values are raw
per-box values. Box pairs were retained at $-$35/$-$10 spacings of 14--22 bp.
The executed per-box significance thresholds differ between the two cached
scans that supply the displayed panels --- 0.025 for the scan behind
\emph{cop} and \emph{ycbr}, and 0.05 for the scan behind \emph{sacaa}, with
the Shine--Dalgarno threshold at the code default of 0.05 in both --- and the
mean information content quoted in each promoter panel title was computed
over a 300 bp scan interval, not over the displayed crop.

ViennaRNA
% TODO cite: ViennaRNA -- the artifact's own note says it should be cited
% only if an analysis using it is actually reported.
was used elsewhere in this project for RNA secondary-structure prediction,
but no base-pair-probability panel appears among the figures shown here, and
none of the displayed panels contains a thermodynamic prediction. The pTA1040
single-strand origin is therefore not presented as a structural or
thermodynamic control: the stem-loops drawn on it are literature dyad
symmetries relocated against the deposited sequence, as described above.
